% =========================================================
% TikZ diagram: Open set definition
% Place after the Open Set / Closed Set definition box.
% Requires: \usepackage{tikz}  \usetikzlibrary{calc}
%           \usepackage{caption}   (for \captionof)
% =========================================================

\begin{center}
\begin{tikzpicture}[
  dot/.style     = {circle, fill, inner sep=1.8pt},
  lbl/.style     = {font=\small},
  sublbl/.style  = {font=\scriptsize, text=gray!55!black}
]

% --- Open set S ---
\fill[blue!10]
  (0, -0.3)
  .. controls ( 1.4,  0.2) and ( 2.2,  1.4) .. ( 1.8,  2.8)
  .. controls ( 1.2,  4.0) and (-0.6,  4.4) .. (-1.8,  3.6)
  .. controls (-3.2,  2.8) and (-3.4,  1.0) .. (-2.4,  0.0)
  .. controls (-1.6, -0.8) and (-0.6, -0.6) .. ( 0,   -0.3)
  -- cycle;

\draw[blue!55!black, dashed, line width=0.9pt]
  (0, -0.3)
  .. controls ( 1.4,  0.2) and ( 2.2,  1.4) .. ( 1.8,  2.8)
  .. controls ( 1.2,  4.0) and (-0.6,  4.4) .. (-1.8,  3.6)
  .. controls (-3.2,  2.8) and (-3.4,  1.0) .. (-2.4,  0.0)
  .. controls (-1.6, -0.8) and (-0.6, -0.6) .. ( 0,   -0.3)
  -- cycle;

\node[lbl, text=blue!60!black, font=\small\itshape] at (-1.5, 0.6) {$S$};
\node[lbl, font=\small\itshape] at (-3.8, 4.2) {$X$};

% --- Interior point x: ball fits ---
\coordinate (x1) at (-0.8, 2.2);
\draw[blue!70!black, dashed, line width=0.7pt] (x1) circle (0.85cm);
\draw[gray!60, line width=0.6pt, ->] (x1) -- ++(0.85, 0);
\node[sublbl, anchor=south] at ($(x1)+(0.42, 0.0)$) {$r$};
\node[dot, fill=blue!70!black] at (x1) {};
\node[lbl, anchor=south east] at ($(x1)+(-0.08, 0.10)$) {$x$};
\node[sublbl, anchor=south, align=center] at ($(x1)+(0, 0.92)$)
  {$B_r(x)\subseteq S$ \quad \textit{fits inside}};

% --- Near-boundary point z: ball leaks ---
\coordinate (z) at (1.35, 1.7);
\draw[red!60!black, dashed, line width=0.7pt] (z) circle (0.52cm);
\node[dot, fill=red!70!black] at (z) {};
\node[lbl, anchor=north west] at ($(z)+(0.08,-0.10)$) {$z$};
\coordinate (leak) at ($(z)+(0.46, 0.24)$);
\draw[red!70, line width=1.3pt]
  ($(leak)+(-0.09,-0.09)$) -- ($(leak)+(0.09, 0.09)$);
\draw[red!70, line width=1.3pt]
  ($(leak)+(-0.09, 0.09)$) -- ($(leak)+(0.09,-0.09)$);
\node[sublbl, anchor=west, align=left] at ($(z)+(0.65, 0.15)$)
  {$z\in S$, but every $B_r(z)$\\[1pt]meets $X\!\setminus\! S$};

% --- Footer ---
\node[lbl, align=center] at (-0.8, -1.15)
  {$S$ is \textbf{open}
   $\iff$
   $\forall x \in S,\;
    \exists r > 0,\;
    B_r(x) \subseteq S$};

\end{tikzpicture}
\captionof{figure}{The open set condition in a metric space $(X,d)$.
Interior point $x$: a ball $B_r(x)$ fits entirely inside $S$, witnessing
openness at $x$.
Near-boundary point $z$: every ball around $z$ leaks into $X\setminus S$,
so if such a point existed in $S$ then $S$ would fail to be open.
The dashed boundary of $S$ itself signals that $S$ is open (boundary excluded).}
\label{fig:open-set}
\end{center}